\chapter{Conclusions on the generative construction and its realizations}
\chaptermark{Conclusions on the generative construction}
\label{hmt:sistematica:conclusiones}

The development establishes a constructive relationship between modular arithmetic, oriented dynamics, value generation, exceptional incidence and information recovery. The unity of this relationship resides in the enriched state produced by APP--TRIT--TPK. Each realization receives a determined part of its structure and preserves the data required by the next operation. Numerical evaluation, geometric form and the record of an execution are thus articulated through explicit maps.

The principal conclusions concern three inseparable mathematical questions: how a value is produced, what structure its production preserves and under which operations it can be recovered. The realizations of Dimensional Mechanics add a fourth determination: how that structure is expressed in observables with declared units and coupling conditions.

\section{Arithmetic restitution and memory of operations}

The initial APP construction fixes two levels: position on the grid and the arithmetic evaluation performed there. Nonadic reduction preserves the integer exactly when it retains remainder and quotient. Balanced ternary reduction has the same property with its oriented representative and extension quotient. The composition identities prove that both registers can be transported between the two decompositions.

In particular, restitution \(n=\rho_9(n)+9q_9(n)\) and compatibility \(c_3(n)=c_3(\rho_9(n))+3q_9(n)\) allow integer information to be tracked through changes of representation. Addition carries and cross-terms of multiplication are computed on those coordinates. Recoverable information thus has, from this stage onward, a determined algebraic content \ref{thm:app-reconstruccion}, \ref{trit:prop:operaciones}.

Composition of trajectories adds ordered memory. If a transition acts through \(U_\varepsilon(m)=C_\varepsilon m+
\kappa_\varepsilon\), concatenation satisfies
\[
 C_{\beta\alpha}=C_\beta C_\alpha,
 \qquad
 \kappa_{\beta\alpha}=C_\beta\kappa_\alpha+
 \kappa_\beta.
\]
The second trajectory transports the accumulated memory before adding its increment. This cocycle law explains why the chronological register contains more information than an aggregated carry. Its proof fixes the operation that must be preserved when composing realizations \ref{x_reciprocidad:iv:cont:concatenacion}.

The nonadic calendar supplies another exact distinction. The return of nine phases increments the turn coordinate. The extension \(C_{12}\to C_{108}\to C_9\), with its carry, expresses compatibility between the visible phase and the dodecaphase register. The enriched continuation also preserves successive turns and incidences. Periodicity of a reading and evolution of its memory therefore constitute two simultaneous properties of the same transport.

Transfer between quadratic and cubic completions supplies another exact conservation identity. In the nonadic case,
\[
 (72,27)\longmapsto(73,27)
 \longmapsto(10\cdot73-1,10\cdot27+1)=(729,271).
\]
The first operation incorporates the registered unit; the second changes the scale and redistributes one unit through a zero-sum correction. The normalized sum remains equal to one. This construction gives precise content to the holographic extreme and mean ratio and its conservative extension across scales \ref{x_reciprocidad:lib03:prop-completacion}.

\section{Determination of the regions and the coinductive fourfold relation}

Finite enumeration identifies the admissible configurations, emissions and words on which the lifts act. The chain \(w_6\to w_{12}\to w_{18}\to w_{24}\to w_{30}\to R_{36}\to G_9\) preserves prefixes and transport conditions when extending the regional realizations. Cylindrical compatibility and contraction provide the passage to arbitrary depth.

The regional outputs \(\pi_{\rm HMT}\), \(\varphi_{\rm HMT}\) and \(e_{\rm HMT}\) combine a positional expansion with an operator provenance. The classification of the five regions of \(\pi\), the incidence of the \(\varphi\) axis, the joint continuations and the mirror involution preserve their distinctions in the limit. Subsequent analytic identifications certify properties of those outputs and allow their comparison in established mathematical languages.

The nonadic connection relates regional generation to return structure. Modal incidence arises from the tritic selector and its sequence of operations. Its powers generate recurrences; regional marks are then applied to the return support. This precedence distinguishes the regional emitter, the recurrence of a modal realization and the evaluation of a boundary ratio. Composition makes their relationships computable without identifying objects of different types.

The register \(K\) is determined from the regional bands. The first zero advance of the conversion clock locates the nonadic antecedent; oriented defects and the phase--charge reader produce the descriptor. The terminal repertoire, ternary cylinders and exceptional incidences delimit the candidates. The heterotypy condition selects the register in that finite domain. Subsequent inversion recovers its order, windows and signed realization. Uniqueness expressly preserves the candidate domain and selector conditions.

The integer fine-structure closure composes regional readings with the rational coordinate of \(K\). Its compatible carry determines a unique infinite section, and the tail identity distinguishes the boundary of an isolated window from the restriction of the limit. In the notation of its development,
\[
 \alpha_{\rm HMT}
 =\pi_{\rm HMT}+e_{\rm HMT}-\varphi_{\rm HMT}
  -4-\kappa_{\rm per}.
\]
The equality preserves the origin of all its terms and the complete decimal normalization \eqref{x_reciprocidad:eq:alpha-salida-a-completa}.

The analytic representation develops that same pre-coordinate through a function defined in all degrees. If \(a\) is the generated pre-coordinate and \(q=1/729\), the linking coefficient
\[
 \lambda=-\frac{E_9(a)(1-qa^3)}{a^{10}}
\]
determines the geometric extension \(E_\infty(X)=E_9(X)+\lambda X^{10}/(1-X^3/729)\). Subsequent coefficients satisfy \(b_{10+3n}=q^n\lambda\) and \(b_{11+3n}=b_{12+3n}=0\). Uniqueness is proved within this class of extensions, and the root is identified with the common coordinate. The second route contributes a correlated analytic representation with convergence control and projective compatibility \ref{x_reciprocidad:subsec:alpha-lector-completo-rev08}.

\section{Archimedean realization and exceptional incidence of the continuum}

The truncation systems preserve admissible histories and their registers. The proved extendibility guarantees the existence of compatible sections. Cylindrical partitions provide an Archimedean evaluation; the incidence structure remains in the conjugate projection. Contraction determines the real value, while orientation, boundary and memory preserve the distinctions of the state.

This double projection jointly organizes spectral transport, solenoidal balance, gauge incidence, cohomological coherence and the determinant line. Naturality identities allow refinements and realizations to be carried out in compatible orders. The discrete continuum thus receives a common operational structure, on which its subsequent evaluations act.

Exceptional incidence is developed through a sequence of concrete constructions. Codes and markings determine lattice gluings; the Leech neighbour is proved positive, even, unimodular and rootless in the declared metric and chamber. Lattice identification uses these properties jointly. The construction of oscillators and the twisted lattice algebra, followed by the orbifold parity sector, determines the realization of \(V^\natural\) and its Monster action~\ref{x_reciprocidad:exc:moonshine}.

The significance of this sequence lies in the related genealogy of value and symmetry and their precise objects of action. Exceptional symmetry acts on the corresponding code, lattice or algebra. Dimensional reductions and duality realizations receive those carriers with their conditions. Numerical refinement and exceptional extension thus retain their linking maps within the common state.

\section{Bilateral conservation and reversibility at arbitrary depth}

The bilateral law supplies an operator conclusion independent of dimension. For two isometries \(\mathcal B,\mathcal C:H\to H'\) and \(a>0\), define \(Q_-=a\mathcal B-\mathcal C\) and \(Q_+=(a+1)\mathcal B+\mathcal C\). The identity
\[
 (a+1)Q_-^*Q_-+aQ_+^*Q_+
 =\bigl((a+1)^3+a^3\bigr)I
\]
follows by cancelling the cross-terms. Its validity includes infinite-dimensional Hilbert spaces and transports between distinct carriers with common domain and codomain.

In the nonadic specialization, \(a=8\), the normalization factors through the weights \(72/73\) and \(1/73\), followed by a rotation of the two channels. The resulting isometry has a left inverse of norm one. Recovery therefore maintains optimal control of the joint error in the declared norm. The factor \(73=72+1\) enters into that identity; its other realizations as a norm or lattice index retain their own proofs.

The transport of observables specifies the content of conservation. For a reader \(G\) and a unitary relative transport \(R\), the induced reading is
\[
 \mathcal T_a(G)=
 \frac{a(a+1)G+R^*GR}{a(a+1)+1}.
\]
The equality \(\mathcal T_a(G)=G\) is equivalent to commutation \(GR=RG\). This yields a relative-symmetry criterion distinguishing invariant observables from those that change during transduction \ref{x_reciprocidad:lib04:teo-lector}.

The recording policy at arbitrary depth preserves the continued component and its complement at each stage. The telescopic balance produces a finite isometry and, under the proved contraction bound, an infinite-record isometry. For \(a=8\), the uniform rate is \(r_8=729/1241<1\). The first \(N\) complements recover the state with truncation error bounded by \(r_8^N\|v\|\); a joint error \(\varepsilon\) adds exactly \(\varepsilon\) to that bound. The complete record admits the convergent inverse of Theorem~\ref{x_reciprocidad:lib04:teo-archivo}.

The evolutionary conservation of unity is therefore expressed through local identities, isometric compositions and recovery in the limit. Modification of form or carrier is compatible with restitution of the state when the channels and incidences specified by the theorem are preserved.

\section{Geometric reciprocity and exterior conservation}

The coupling register produces a direction orthogonal to the uniform mode. Its positive diagonal flow has determinant one. The radii vary compensatorily; Plücker minors preserve sign and vanishing, and balanced quotients preserve their value. The same direction generates a unit-covolume deformation on the lattice planes.

The reciprocity of this deformation is expressed by \(\Lambda_t^*=\Lambda_{-t}\). The commuting ternary symmetry persists along the flow, while integrality and exceptional realization retain the specific domain in which they were proved. Theorem~\ref{x_reciprocidad:rec:thm:red} determines these properties simultaneously.

In the chart \(s=(1-z)^{-1}\), radial inversion \(z\mapsto1/\overline z\) becomes \(s\mapsto1-\overline s\). The unit circle corresponds to the line \(\Re s=1/2\), with the singular point excluded. Theta--Mellin reciprocity, polar structure and the Poincaré and Klein realizations preserve the analytic and geometric conditions of their own readers. Recorded orientation allows the radial reading to be inverted wherever the scalar reading combines conjugate realizations.

These identities broaden the meaning of the coupling constant: the arithmetic register determines a recoverable frame, a deformation direction and transports preserving exterior invariants. The correspondence between realizations is obtained by composing the constructed maps.

\section{Action, angular geometry and material observables}

The Dimensional Mechanics realization preserves an explicit precedence. The action scale receives the fine-structure coordinate and the regional reading of \(\varphi\); the local norm of order sixteen fixes its decimal valuation. The forward and return Planck sections preserve their quotient \(R_{\rm act}\). The angular pair and its conjugate channels are then evaluated, with their orientation and units~\ref{ii:sec:ii-definicion-planck}.

The marked inclusion of a hexad in an octad produces incidences of cardinalities \(90\) and \(120\). The dodecaphase chain determines the multiplicities of the Barbero--Immirzi functional and the pole coefficient \(1/24\). Its evaluation uses the angle, counter-angle and previously constructed channels. Parity under angular inversion preserves the orientation of that reading~\ref{ii:sec:barbero}.

Angular transport toward Witt incidence supplies an explicit composition with the mixing reader. Its left inverse and Gram identity preserve the transported angular information \ref{ii:thm:ii-witt-angular-transport}. The electric and magnetic constitutive structure likewise receives its modes, return ratios and scales; its identities allow the angular coordinates to be reconstructed in the specified domain.

The contragredient gravitational section of action satisfies
\[
 G_a=\frac{c_{\rm int}^{3}R_{108}^{2}}{\hbar_a},
 \qquad
 \frac{G_a\hbar_a}{c_{\rm int}^{3}}=R_{108}^{2}.
\]
The coincident-radii condition fixes this section within the declared class. Inversion of the action ratio compensates the gravitational ratio and preserves areal unity \ref{ii:sec:ii-definicion-gravedad}. The Boltzmann definition then relates energy per decision to the thermal section. For the clock of one hundred and eight events, \(k_B\Theta_{{\rm clk},a}=h_a/(108t_0\log3)\). Individualization of the transducer preserves the thermal section and unit bases it uses \ref{ii:sec:ii-definicion-boltzmann}.

Material observables arise from the persistence of extensions and their route records. The electronic centre has a point selection determined by the atlas; the other classes use the operators of their fibres, orbits or multisections. The integer signature, character and towers transmit this information to the mass realization. Metrological comparison is carried out on the specified evaluations and couplings, preserving the distinction between operator, spectrum and individualized scalar.

\section{Reconstruction and cardinal classification of the generated continuum}

The cylinder construction proves coverage of values with recovery of registers. Cardinal closure adds a classification of those values in the genealogical closure \(\mathfrak G\) determined by its term semantics and definability. Internal countability of the levels and the condensation lemma assign to each rational cut its first stage \(\beta(r)\) and its first index \(n(r)\) in that stage's enumeration.

The map
\[
 j_{\rm HMT}(r)=\omega\beta(r)+n(r)
\]
is injective on the domain of the generated line \(\mathbb R_{\rm HMT}^{\mathfrak G}\). Encoding countable well-orders constructs the injection of \(\omega_1^{\mathfrak G}\) into that same line. Cantor--Bernstein concludes
\[
 \left|\mathbb R_{\rm HMT}^{\mathfrak G}\right|
 =\aleph_1^{\mathfrak G},
 \qquad
 \mathfrak G_{\rm HMT}(h,m_\infty)
 \models\mathrm{CH}.
\]
This is the affirmative closure of the continuum hypothesis established in Theorem~\ref{ix:thm:decision-continuo-hmt}.

The conclusion encompasses every internal subset of that line: each is countable or has the cardinality of the generated continuum. The representation \(\mathfrak G_{\rm HMT}(h,m_\infty)=L[h,m_\infty]\) identifies the construction's domain after developing its genealogy. Transfers to a different set-theoretic domain retain as hypotheses the coverage and coincidence of the objects on which the injections act. The formulation thus keeps the cardinal result, its semantics and its quantifiers together.

The spectral, solenoidal, gauge, cohomological and determinantal terminal realizations are articulated through the same discrete structure. Their theorems specify the operator, positivity, coercivity or descent used by each realization, together with the corresponding naturality identities. The joint factorization of Chapter~\ref{chap:terminal-arquitectura-natural-comun} preserves that origin and the domain of each conclusion.

\Needspace{27\baselineskip}
\section{Demonstrative content and scope of the integration}

The demonstrative strength of the work arises from composition between results of different types. Symbolic identities prove equalities for their quantified domains. Exhaustive enumerations resolve finite selectors and censuses. Rational bounds control approximations and roots. Compatibility, convergence and naturality arguments justify extensions and passages to the limit.

Lean formalizations verify the specific declarations they contain, with their types, premises and dependencies. Executable calculations certify the instances or finite domains defined by their algorithms. Correspondence between mathematical statement and certificate assigns each check its precise scope. The physical realization also preserves the dimensional basis and the coupling rule with the observable.

The overall result is an operational genealogy: discrete arithmetic produces oriented states; their extensions determine constants and a continuum with incidence; their registers admit conservative transformations and recovery; their dimensional realizations provide structured observables. The unity of the development is expressed in these demonstrated relationships. Value, symmetry and recoverable information are linked by the construction producing them.
